\section{Demand, composition, and an information limit}
\label{app:limits}

\subsection{Demand and composition do not disappear}

At common scaled demand $g(t)$, variable labor is proportional to $g(t)h(t)$. Section~\ref{sec:workforce} gives the consequence: a demand response is not a failure of gate-level
substitution. New governing
requirements can add contracts, while recurring requirements necessary for an existing contract remain
in its support account. The fixed 2033 workforce deadline prevents new terminology from resetting the dated
prediction; new DGF roles remain in its accounting boundary. A changing economy prevents an old task
registry from proving a claim about all future jobs.

\subsection{Why writing inputs and outputs is not an existence proof}

Suppose two possible environments produce the same accessible input history but require disjoint
correct outputs. Any deterministic system restricted to that history returns the same output in both
and must fail in at least one. A randomized system cannot guarantee success in both either. Additional
inquiry or a contract that permits an explicit unresolved disposition is necessary. A human with access
to new information changes the information boundary, not this logical constraint.

This example identifies a concrete test of the law: can the missing observation or accepted decision
procedure be supplied without indispensable human execution? The paper predicts eventual coverage of
the DGF functions; it does not infer coverage from input-output notation alone.
